\section{The same original relative canonical section is a unit along original local lifts}
\label{sec:valuative-canonical-section}

\begin{theorem}[Actual canonical section unitness on the original local spectrum]
\label{sc0-valuative-canonical-section:original-unit}
\lean{EnvelopeProperty.Char0.RoundedCanonicalModule.pullback_differentialSection_isIso_of_source_localization}
\leanok
For EVERY actual original lift over an original source localization
as in the preceding theorem, the SAME original relative canonical
differential SECTION pulls back, with the genuine native unit
comparison, to an actual isomorphism from the original structure
module. All original source and target smoothness, actual source
differential basis and original localization data are retained.
No section-unit, inverse, frame or Hom-section pullback-matching
certificate is supplied.
\end{theorem}
\begin{proof}
The preceding theorem derives that the SAME ENTIRE original
canonical differential map becomes an isomorphism after this
actual pullback. The genuine Hom-section isomorphism theorem then
uses the proved full original monoidal base-change identity and
actual right-dual evaluation comparison to derive that pullback
of the SAME original Hom-section is an isomorphism. By its actual
definition this Hom-section is precisely the original relative
canonical differential section. Compose with the inverse of the
TRUE native structure-module pullback unit comparison, which is
itself an isomorphism. This gives the asserted same entire unit
section map on the original local spectrum. In particular the
original intervening scheme maps are not assumed smooth or flat,
and no local Jacobian or coordinate comparison is input.
\uses{sc0-valuative-canonical:whole-pulled-differential,
      sc0-hom-section-iso:whole-map}
\end{proof}

This actual unit section canonically trivializes the SAME pulled
relative canonical line. It will define the regular coefficient
of every actual pulled original section in the original valuation
ring. Comparing it with the SAME original prescribed rational
coefficient and the full rounded-module pushforward remain next.
No vanishing conclusion is asserted.
